\documentclass[11pt,a4paper]{article}
\usepackage[T1]{fontenc}
\usepackage{lmodern,amsmath,amssymb,amsthm,mathtools,booktabs,microtype}
\usepackage[margin=27mm]{geometry}
\usepackage[colorlinks=true,linkcolor=blue,citecolor=blue,urlcolor=blue]{hyperref}
\newtheorem{theorem}{Theorem}[section]
\newtheorem{proposition}[theorem]{Proposition}
\newtheorem{lemma}[theorem]{Lemma}
\newtheorem{corollary}[theorem]{Corollary}
\newcommand{\Q}{\mathbb Q}\newcommand{\Z}{\mathbb Z}
\DeclareMathOperator{\Tr}{Tr}
\DeclareMathOperator{\Odd}{Odd}
\setlength{\emergencystretch}{2em}
\title{Quadratic Covers of Two Contact Functions}
\author{}\date{10 November 2026}
\begin{document}\maketitle
\begin{abstract}
A quadratic equation can force two contact values to have the
same valuation parity, including on the prescribed eighth-power
scales. We determine the resulting curves: for every rational
scale greater than four, the product cover has genus four.
Consequently only finitely many rational slopes can have equal
odd valuation sets outside a fixed finite set of primes.
More generally, a finite curve cover on which the two contact
squareclasses become dependent cannot have genus zero or one.
The negative product cover has an explicit rational section for
every scale, but an exact inert-prime calculation shows that this
section is incompatible with the elliptic comparison tests for
infinitely many odd input shifts. These statements distinguish
the valid parity implication from the separate existence of all
required arithmetic witnesses.
\end{abstract}

\section{Contacts and valuation parity}

Put
\[
 \Theta=\left\{0,-1,-\frac12,-\frac13,-\frac14,-\frac25\right\},
 \qquad \Theta^*=\Theta\setminus\{0\},
\]
\begin{equation}\label{eq:R}
 F(t)=\prod_{\theta\in\Theta}(t-\theta)
 =\frac{t(t+1)(2t+1)(3t+1)(4t+1)(5t+2)}{120},
 \qquad R(t)=\frac{F(t)}{(t-1)^8}.
\end{equation}
Fix a rational number $a>4$. A rational slope $\lambda$ is
admissible if $\lambda$ and $a\lambda$ lie outside
$\Theta\cup\{1\}$. In particular all contact values considered
below are defined and nonzero.

Let $S$ be a finite set of primes containing $2,3,5$. For
$q\in\Q^*$ define
\[
 \Odd_S(q)=\prod_{\substack{p\notin S\\v_p(q)\ \mathrm{odd}}}p.
\]
This definition includes negative odd valuations. For
$t\in\Q\setminus(\Theta\cup\{1\})$, put
\[
 M_S(t)=\prod_{\substack{p\notin S\\
    v_p(t-\theta)>0\ \mathrm{odd}\ \text{for some }\theta\in\Theta}}p.
\]

\begin{lemma}\label{lem:odd}
For every such $t$, one has $M_S(t)=\Odd_S(R(t))$.
Consequently an equation
\begin{equation}\label{eq:covercondition}
                       z^2=-R(\lambda)R(a\lambda)
                       \qquad(z\in\Q^*)
\end{equation}
implies $M_S(\lambda)=M_S(a\lambda)$.
\end{lemma}
\begin{proof}
Outside $2,3,5$, all contacts are integral and all their pairwise
differences are units. If $v_p(t)<0$, every factor of $F(t)$
has valuation $v_p(t)$ and $v_p(t-1)=v_p(t)$. Hence
$v_p(R(t))=-2v_p(t)$ is even. If $t$ is integral, at most one
contact factor has positive valuation. The denominator of $R$
has even valuation, so the parity of $v_p(R(t))$ is exactly
the parity of that contact valuation. This also handles primes
at which a contact and the pole become congruent: no assertion
that $v_p(R(t))$ is positive is needed here.

Equation~\eqref{eq:covercondition} makes the sum of the two
valuations even at every prime. Their odd valuation sets are
therefore equal, and the first assertion proves the conclusion.
\end{proof}

\section{The genera of the covers}

For $d\in\Q^*$, let $C_{a,d}$ be the smooth projective curve
with function field obtained from $\Q(\lambda)$ by adjoining
a square root of $dR(\lambda)R(a\lambda)$.

\begin{theorem}\label{thm:genus}
For every rational $a>4$ and $d\ne0$, the curve $C_{a,d}$ is
geometrically integral of genus four. A birational affine model is
\begin{equation}\label{eq:model}
 Y^2=d\prod_{\theta\in\Theta^*}
                      (\lambda-\theta)(\lambda-\theta/a).
\end{equation}
The quadratic covers associated with either $dR(\lambda)$ or
$dR(a\lambda)$ have genus two.
\end{theorem}
\begin{proof}
Factoring out the common zero at zero gives
\[
 R(\lambda)R(a\lambda)
 =\frac{a^6\lambda^2
       \prod_{\theta\in\Theta^*}(\lambda-\theta)(\lambda-\theta/a)}
       {(\lambda-1)^8(a\lambda-1)^8}.
\]
If $z^2=dR(\lambda)R(a\lambda)$, the change
\[
             Y=\frac{z(\lambda-1)^4(a\lambda-1)^4}{a^3\lambda}
\]
proves~\eqref{eq:model} on the admissible open set.

All ten displayed roots are distinct. Within either group this
is immediate. An equality $\theta=\eta/a$, with
$\theta,\eta\in\Theta^*$, would give $a=\eta/\theta$.
The absolute values of the five nonzero contacts lie in
$[1/4,1]$, so this ratio is at most four. Thus the degree-ten
polynomial in~\eqref{eq:model} is squarefree.

The double cover of the projective line is ramified at precisely
these ten simple roots. Its even-degree pole at infinity is
unramified. Riemann--Hurwitz~\cite[Section~53.12]{Stacks} gives
\[
                         2g-2=2(-2)+10=6.
\]
The simple roots also show that the defining rational function
is not a square over $\overline\Q(\lambda)$, which proves
geometric integrality. For a single contact function there are
six simple zeros, and all pole orders are even. The same formula
gives $2g-2=-4+6=2$, hence genus two.
\end{proof}

\begin{theorem}\label{thm:finite}
Fix $a\in\Q$, $a>4$, and a finite set of primes $S$.
There are only finitely many admissible rational slopes satisfying
\begin{equation}\label{eq:equalsets}
                          M_S(\lambda)=M_S(a\lambda).
\end{equation}
When $S$ does not contain $2,3,5$, the same conclusion holds
with the evident contact definition of $M_S$.
\end{theorem}
\begin{proof}
Enlarging $S$ weakens~\eqref{eq:equalsets}, so we may assume
that it contains $2,3,5$. By Lemma~\ref{lem:odd}, the product
$R(\lambda)R(a\lambda)$ has even valuation at every prime
outside $S$. Therefore
\[
                     R(\lambda)R(a\lambda)=d w^2
\]
for $w\in\Q^*$ and one of the finitely many numbers
\[
                  d\in\left\{\pm\prod_{p\in S}p^{e_p}:
                                           e_p\in\{0,1\}\right\}.
\]
Negative valuations cause no difficulty: only their parity
enters the rational squareclass. This equation gives a rational
point on $C_{a,d}$ by taking $z=dw$.
Each of these finitely many curves has genus four by
Theorem~\ref{thm:genus}, so its rational points form a finite set
by Faltings' theorem~\cite{Faltings}. Their slope projections
form a finite set as well.
\end{proof}

Thus equality of the two odd valuation sets cannot coexist with
an unbounded supply of slopes at a fixed scale, even if arbitrary
auxiliary algebraic variables are used to impose that equality.
This conclusion does not assert that the finite set contains
no arithmetically suitable slope.

\begin{proposition}\label{prop:finitecovers}
Let $C/\Q$ be a smooth projective geometrically integral curve
with a nonconstant rational function $\lambda$. Suppose that,
in $\Q(C)$,
\[
              R(\lambda)^{e_0}R(a\lambda)^{e_1}=d w^2,
 \qquad d\in\Q^*,\quad w\in\Q(C)^*,
\]
for some nonzero $(e_0,e_1)\in\{0,1\}^2$.
Then $g(C)\ge2$. If $e_0=e_1=1$, then $g(C)\ge4$.
In particular, such a dependence of squareclasses cannot be
imposed by a curve of genus zero or one dominating the slope line.
\end{proposition}
\begin{proof}
The rational function on the left is not a square over
$\overline\Q(\lambda)$, as shown by its simple branch points.
The displayed equation therefore gives an embedding of the
corresponding quadratic function field into $\Q(C)$, and hence
a nonconstant morphism from $C$ to its smooth projective curve.
The target has genus two if just one exponent is nonzero, and
genus four if both are nonzero. Riemann--Hurwitz gives
$g(C)-1\ge\deg(f)(g(\mathrm{target})-1)$, proving both assertions.
\end{proof}

\section{The rational section and an inert prime}

The negative product cover nevertheless has a rational point
for every scale. Set
\begin{equation}\label{eq:section}
                         \lambda_a=-\frac{a+1}{5a}.
\end{equation}
Then $-1/4<\lambda_a<-1/5$ and $a\lambda_a<-1$, so the pair
is admissible and
\begin{equation}\label{eq:signs}
                         R(\lambda_a)<0<R(a\lambda_a).
\end{equation}
The exact identities are
\begin{align}
 R(\lambda_a)R(a\lambda_a)
 &=-\left(\frac{25a(a-4)(a-1)(a+1)(2a-3)(3a-2)(4a-1)}
 {24(a+6)^4(6a+1)^4}\right)^2,\label{eq:sectionidentity}\\
 R(a\lambda_a)
 &=\frac{25(a+1)(a-4)(2a-3)(3a-2)(4a-1)(a-1)}
 {24(a+6)^8}.\label{eq:sectionR}
\end{align}
For a direct proof of the first identity, the involution
$\psi(t)=-t/(5t+1)$ satisfies
$F(\psi(t))=-F(t)/(5t+1)^6$ and
$a\lambda_a=\psi(\lambda_a)$. Substitution gives
\eqref{eq:sectionidentity}; direct substitution gives
\eqref{eq:sectionR}. These identities also appear in~\cite{Product}.

Use the elliptic curve and non-torsion point
\[
 K=\Q(\sqrt2),\quad D=11-6\sqrt2,\quad E:y^2=x^3-Dx,
\]
\[
 P=\left(\frac{57-28\sqrt2}{8},
          \frac{-660+417\sqrt2}{32}\right),\qquad
 T(m)=\Tr_{K/\Q}x([m]P).
\]
The non-torsion assertion has an unconditional finite-reduction
certificate in~\cite{CM}. At the places $\sqrt2\equiv6\pmod{17}$
and $\sqrt2\equiv5\pmod{23}$, the reductions of $P$ have orders
five and three, while the curve orders are twenty and twenty-four.
If $P$ were torsion, its order would be divisible by five; its
five-primary part would inject at the second place, contradicting
the order twenty-four. Write $f(Q)=1/x(Q)$, with its regular
value zero at the identity, and let $\sigma$ denote conjugation of $K$.

\begin{lemma}\label{lem:inert}
Let $p>76176$ be inert in $K$, and let $H\in E(K)$.
The two conditions, when their values are defined,
\[
                  v(f(H))>0,\qquad
                  v(\sigma f(H-P))>0
\]
cannot both hold at the unique place above $p$.
\end{lemma}
\begin{proof}
The curve has good reduction, and a point with positive
reciprocal-abscissa valuation reduces to the identity. The
unique place above an inert prime is preserved by $\sigma$.
The displayed conditions would therefore make both $H$ and
$H-P$ reduce to the identity, and hence would make $P$ reduce
to the identity. But the displayed affine coordinates of $P$
are integral at every odd prime; on the good projective model
their reduction is a finite point. This is a contradiction.
\end{proof}

\begin{theorem}\label{thm:sectionfailure}
Put
\[
        p=315467,\qquad m=p+1=315468,\qquad
        M=2p^2m=62790389880761304.
\]
For every integer $k\ge0$, take the true integer input $X=0$
and the odd shift
\[
                  B=3+kM,\qquad n=\beta=B,
 \qquad a=B^8+T(B)^8.
\]
At the slope~\eqref{eq:section},
\[
                v_p(R(\lambda_a))=v_p(R(a\lambda_a))=1.
\]
The elliptic comparison tests obtained from a single retained
pole certificate for $R(a\lambda_a)$ cannot be satisfied.
\end{theorem}
\begin{proof}
Trial division proves that $p$ is prime, and $p\equiv3\pmod8$,
so it is inert in $K$. Exact addition on $E$ gives
\begin{equation}\label{eq:T3}
 T(3)=\frac{59609817879309699127684977}{508847420105479203017476}.
\end{equation}
The point $[3]P$ is affine and integral at $p$.
Writing $a_3=3^8+T(3)^8$, exact rational arithmetic gives
\begin{equation}\label{eq:modp2}
              a_3\equiv4+67222p\pmod{p^2},
              \qquad p^2=99519428089.
\end{equation}
In particular $v_p(a_3-4)=1$.
The finite-field addition certificate also gives
$[m]P=O$ after reduction modulo $p$. It computes the order
of this reduction to be $157734$; the larger annihilator $m$
is sufficient here.

For completeness, work over the unramified quadratic extension
of $\Q_p$. Let $E_1$ and $E_2$ be the kernels of reduction
modulo $p$ and $p^2$. The formal parameter at the identity
identifies the group law on $E_1/E_2$ with addition in
$p\mathcal O/p^2\mathcal O$. Hence $[p]E_1\subset E_2$.
As $[m]P\in E_1$, the congruence $B\equiv3\pmod{2p^2m}$
implies
\[
                       [B]P\equiv[3]P\pmod{p^2}.
\]
Since the latter point is affine, its abscissa and its trace
are preserved modulo $p^2$. Also $B\equiv3\pmod{p^2}$.
Consequently
\[
                 B^8+T(B)^8\equiv a_3\pmod{p^2},
                 \qquad v_p(a-4)=1.
\]
All other numerator factors and the denominator in
\eqref{eq:sectionR} are units at this prime: their residues
at $a=4$ involve only $2,3,5$. Thus $v_p(R(a\lambda_a))=1$.
Likewise
\[
 \lambda_a+\tfrac14=\frac{a-4}{20a},\qquad
 a\lambda_a+1=-\frac{a-4}{5}
\]
both have valuation one, and the remaining contact factors are
units. This gives $v_p(R(\lambda_a))=1$.

At a positive odd contact, the retained trace normalization
and pole obstruction force every comparison test to have
positive valuation at one place above $p$, exactly as in
\cite[Lemma~4.1]{Conditional}. Two of those tests are
$f(H)$ and $\sigma f(H-P)$. Lemma~\ref{lem:inert} makes
their simultaneous positivity impossible. Hence no comparison
witness exists at the displayed slope.
\end{proof}

The theorem concerns the explicit section~\eqref{eq:section}.
It does not exclude every rational point of $C_{a,-1}$.
The shifts in the theorem are unbounded, so increasing a
uniform lower bound for the odd shift does not make this
particular section sufficient for all integer inputs.

\section{The implication supplied by one pole}

We record separately what the quadratic equation does prove.
Retain the hypotheses and guarded elliptic point data of
\cite{Conditional}: the weak rational $abc$ inequality with
exponent $121/100$, the group identity $2E(K)=\Z P$, and the
specified pole obstruction. Put
\[
 A=[n]P,\quad H=[h]P,\quad U=[u]P,\quad
 V_0=H-P,\quad V_1=U-A,
\]
with the nonzero charts and quotients required there. For
$z=-x/y$, put
\[
 L=\frac{z(U)}{z(H)},\qquad
 L'=\sigma\!\left(\frac{z(V_1)}{z(V_0)}\right),
\]
and use the six tests
\[
 D_1=f(H),\ D_2=f(U),\ D_3=\sigma f(V_0),\
 D_4=\sigma f(V_1),\ D_5=L-L',\ D_6=L-\beta.
\]
Set $\beta=2X^2+B$, $a=\beta^8+T(n)^8$, and retain
only the pole certificate for
\[
 W_1=\prod_{j=1}^6\left(1+
        \frac{D_j^2}{(a\lambda-1)^4R(a\lambda)}\right),
 \qquad b_1=\frac{R(a\lambda)}{\Tr(W_1^{-1})^2}.
\]
Adjoin equation~\eqref{eq:covercondition}, with all the displayed
denominators and contact values guarded nonzero.

\begin{proposition}\label{prop:soundness}
Under these hypotheses, there is a constant $B_*$ such that
every rational solution of this system for an odd $B\ge B_*$
has $X\in\Z$.
\end{proposition}
\begin{proof}
Take $S=\{p:p\le76176\}$ and suppose $\beta\ne n$.
Let $N$ be the positive reduced numerator of $\beta-n$.
The unchanged local comparison argument at the retained pole
gives $M_S(a\lambda)\mid N$.
Lemma~\ref{lem:odd} supplies $M_S(\lambda)=M_S(a\lambda)$,
so both radicals divide $N$.

The two height consequences of weak $abc$ proved in
\cite[Lemmas~2.1 and~2.2]{Conditional} therefore give, for
$\eta=h(1:\beta:T(n))$,
\[
                  \eta\le q\log N+C,
                  \qquad q=\frac{14641}{16169}<1.
\]
The constant $C$ is independent of the input shift and all
witnesses. Since
$N\le e^{h(\beta)}(1+|n|)$ and $h(\beta)\le\eta$,
\[
          \eta\le\frac{q}{1-q}\log(1+|n|)+\frac{C}{1-q}.
\]
The retained trace-height theorem gives
$n^2\le A_0h(T(n))+J\le A_0\eta+J$, with fixed constants.
These inequalities bound $n$, and then $h(\beta)$, independently
of $B$. But $\beta=2X^2+B\ge B$, so a sufficiently large odd
$B$ is impossible unless $\beta=n$. Finally, if $X=u/v$ is
reduced and $\beta\in\Z$, then $v^2\mid2$ and therefore $v=1$.
\end{proof}

No converse is asserted in Proposition~\ref{prop:soundness}.
The quadratic condition supplies the needed parity implication,
but a definition also requires suitable rational points on the
cover together with all elliptic and pole witnesses for every
integer input. Theorem~\ref{thm:finite} prevents retaining the
usual unbounded choice of slopes at each fixed scale, and
Theorem~\ref{thm:sectionfailure} excludes the displayed section
as a uniform substitute. Neither statement excludes a different
arithmetic existence argument.

\section{Exact verification}

The standard-library program
\nolinkurl{computation/pole_parameters/check_quadratic_contact_covers.py}
verifies both section identities by exact polynomial arithmetic.
It enumerates the possible collision ratios of nonzero contacts;
their largest value is four. The genus calculation then follows
from the ten distinct branch points, not from sampled scales.

The program reconstructs~\eqref{eq:T3} independently by
chord-and-tangent addition and by the third division polynomial.
It verifies primality and inertness of $315467$, all affine
denominator conditions, the exact residues modulo $p^2$, and
the two contact valuations. Two binary multiplication orders
verify the finite-field annihilator used in the progression
argument. The output is
\nolinkurl{computation/pole_parameters/quadratic-contact-covers-certificate.json}.
The proof of finiteness uses Faltings' theorem separately;
no finite computation is claimed to prove that theorem or
the completeness of a defining polynomial.

\begin{thebibliography}{9}
\bibitem{Conditional}
\emph{Thirteen Variables from a Weak $abc$ Inequality},
companion paper, 9 November 2026,
\nolinkurl{papers/conditional_few_variables/conditional13_20261109_v16.tex}.
\bibitem{Product}
\emph{Odd Radicals and Products of Six Linear Factors},
companion paper, 10 November 2026,
\nolinkurl{papers/pole_parameters/product_height_obstructions_20261110_v19.tex}.
\bibitem{CM}
\emph{Complex Multiplication and Trace Cancellation},
companion paper, 10 November 2026,
\nolinkurl{papers/number_fields/complex_multiplication_and_trace_cancellation_20261110_v19.tex}.
\bibitem{Stacks}
The Stacks Project Authors, \emph{The Stacks Project},
Section~53.12, Riemann--Hurwitz, Tag~0C1B.
\url{https://stacks.math.columbia.edu/tag/0C1B}.
\bibitem{Faltings}
G.~Faltings, \emph{Endlichkeitss\"atze f\"ur abelsche Variet\"aten
\"uber Zahlk\"orpern}, Inventiones Mathematicae \textbf{73} (1983),
349--366, \url{https://doi.org/10.1007/BF01388432};
erratum, \textbf{75} (1984), 381,
\url{https://doi.org/10.1007/BF01388572}.
\end{thebibliography}
\end{document}
